\documentclass[10pt,a4paper]{amsart}
\usepackage[margin=19mm]{geometry}
\usepackage{amsmath,amssymb,mathtools}
\usepackage[scr=rsfs,cal=euler]{mathalfa}
\usepackage{xfrac}
\usepackage[unicode,hidelinks,bookmarks=false]{hyperref}
\newcommand{\ie}{i.e.\ }
\newcommand{\cf}{cf.\ }
\newcommand{\resp}{respectively\ }
\newcommand{\Subsection}{\subsection}
\providecommand{\nobreakdash}{\nobreak\hbox{-}}
\setlength{\emergencystretch}{2em}
\multlinegap0pt
\usepackage{bm}
\usepackage{bbm}
\usepackage{dsfont}
\usepackage{xspace}
\usepackage{cancel}
\usepackage{mathtools}
\usepackage{amsthm}
\makeatletter
\@ifundefined{thm}{\newtheorem{thm}{Thm}}{}
\@ifundefined{lemma}{\newtheorem{lemma}[thm]{Lemma}}{}
\@ifundefined{proposition}{\newtheorem{proposition}[thm]{Proposition}}{}
\makeatother
\title[Quantum Mixing and Benjamini-Schramm Convergence of Hyperbol]{Quantum Mixing and Benjamini-Schramm Convergence of Hyperbolic Surfaces}
\author{Kai Hippi}
\date{Source: arXiv:2512.15504v3 (CC BY 4.0)}
\begin{document}
\maketitle

\noindent\emph{Source and licence.} Quantum Mixing and Benjamini-Schramm Convergence of Hyperbolic Surfaces. Version arXiv:2512.15504v3 (CC BY 4.0), distributed under the Creative Commons Attribution 4.0 International licence, as recorded in the arXiv licence metadata for this exact version and shown on the abstract page https://arxiv.org/abs/2512.15504v3. Full rights notice: licenses/arxiv-2512.15504-CC-BY-4.0.txt.\par\smallskip

\noindent\textit{Editorial scope.} Counted unit: the Proposition whose source label is \texttt{Prp: Geometric data} in the article (source0.tex lines 1073--1096, source label \texttt{Prp: Geometric data}), delivered together with the article's own complete original proof of it (source0.tex lines 1100--1218, one \texttt{proof} environment of 119 lines). No mathematics is written, repaired or re-derived: statement and proof are the article's own text line for line. The proof is detached from the statement in the source: the article states the result at lines 1073--1096 and prints its proof at lines 1100--1218, and the proof environment's own header names the statement, so the association is the article's own. Required context retained uncounted: Lmm: Eliminate sum 1 (lines 1222--1232); Lmm: Eliminate sum 2 (lines 1262--1275); Lmm: Large injectivity radius (lines 1323--1345); Lmm: Small injectivity radius (lines 1399--1407); Lmm: Weight integral (lines 1460--1472); Lmm: Points do not matter (lines 1545--1557). Every definition, display and label of those lines is retained unchanged. Omitted from this package and not redistributed: 1--1072, 1097--1099, 1219--1221, 1233--1261, 1276--1322, 1346--1398, 1408--1459, 1473--1544, 1558--2705. Nothing inside a retained excerpt is omitted or altered: every retained body is delivered exactly as the article's lines are written, so no omission receipt of any kind accompanies this package. Citations: the retained excerpts carry no citation command at all, so no bibliography entry of the article is transcribed and no reference is redistributed. Reference adaptation: none was needed and none was made; every reference inside the delivered unit resolves inside it, so no unresolved reference or citation is produced. Printed slips kept literal: no printed word, symbol, hypothesis or number of the counted statement or of its proof was altered. Layout and class: the article's own class is \texttt{amsart} with packages including a4wide, fontenc, inputenc, soul, amsmath, amsfonts, amssymb, mathtools, amsthm, xcolor, graphicx, hyperref, standalone, esint; the wrapper is the lane's pinned \texttt{amsart} editorial wrapper at 10pt on A4, which restates the theorem-like environments the retained text needs and adds the article's own macro definitions that the retained text needs (none), with extra wrapper packages (bm, bbm, dsfont, xspace, cancel, mathtools). The delivered numbering is the unit's own. Assets: no retained span contains a figure environment, a graphics-inclusion command, a PDF-page inclusion, a table or a TikZ picture; no image file is copied into this package, the package declares no asset dependency and no image file is redistributed with it. Licence: version arXiv:2512.15504v3 of this article is distributed under the Creative Commons Attribution 4.0 International licence, as recorded in the arXiv licence metadata for this exact version and shown on the abstract page https://arxiv.org/abs/2512.15504v3. A full rights notice is delivered at licenses/arxiv-2512.15504-CC-BY-4.0.txt. Characters: every retained excerpt is pure ASCII, so the delivered engine, XeLaTeX with its default Latin Modern text font, reports no missing character.
\begin{lemma}[]
	\label{Lmm: Eliminate sum 1}
	Let 
	$$ K^{\Gamma}( x,y ) = \sum_{\gamma \in \Gamma}^{} K( x, \gamma y ), \qquad L^{\Gamma}( x,y ) = \sum_{\gamma \in \Gamma}^{} L( x, \gamma y ) $$ 
	be integral kernels on a compact connected hyperbolic surface $ X= \Gamma \setminus \mathbb{H} $, where $ K $ and $ L $ are integral kernels on $ \mathbb{H} $ that depend only on distance $ d( x,y ) $. Then
	\begin{align*}
		&
		\int_{D} K^{\Gamma}( x,z ) L^{\Gamma}( x,z' ) dx = \int_{\mathbb{H}}  \sum_{\gamma \in \Gamma}^{} K(  x, z ) L( x, \gamma z' ) dx,
	\end{align*}
	where $ D $ is a fundamental domain of $ X $.
\end{lemma}

\begin{lemma}[]
	\label{Lmm: Eliminate sum 2}
	Let 
	$$ K^{\Gamma}( x,y ) = \sum_{\gamma \in \Gamma}^{} K( x, \gamma y ), \qquad L^{\Gamma}( x,y ) = \sum_{\gamma \in \Gamma}^{} L( x, \gamma y ) $$ 
	be integral kernels on a compact connected hyperbolic surface $ X= \Gamma \setminus \mathbb{H} $, where $ K $ and $ L $ are integral kernels on $ \mathbb{H} $ that depend only on distance $ d( x,y ) $. Let $ a $ be a function on $ X $. Then
	\begin{align*}
		&
		\int_{D}  a( z' )^{*} \left( \int_{\mathbb{H}} \sum_{\gamma \in \Gamma}^{}K( x,z ) L( x, \gamma z' ) dx \right)^{2}  dz'
		\\
		& = \
		\int_{\mathbb{H}} a( z' )^{*}  \left( \int_{\mathbb{H}} K( x,z ) L( x, z' ) dx \right) \sum_{\gamma \in \Gamma}^{}  \left( \int_{\mathbb{H}} K( y,\gamma z ) L( y, z' ) dy \right) dz',
	\end{align*}
	where $ D $ is a fundamental domain of $ X $.
\end{lemma}

\begin{lemma}[]
	\label{Lmm: Large injectivity radius}
	Let $ X = \Gamma \setminus \mathbb{H}$ be a compact connected hyperbolic surface with fundamental domain $ D $. Let $ D( 4T ) $ denote the points with injectivity radius larger than $ 4T $. Let $ a $ be a mean zero function on $ X $. Then
	\begin{align}
		&
		\nonumber
		\left| \int_{D( 4T )} a( z ) \int_{ B( z, 2T ) } a( z' )^{*}   F_{t,t', d( z,z' )}  \sum_{\gamma \in \Gamma}^{} F_{t,t', d( \gamma z, z' )} dz' dz \right|
		\\
		& = \
		\label{Expr: Large injectivity radius}
		30 \| a \|_{2}^{2} \, \int_{0}^{2T}  \sinh( \rho ) \, F_{t,t',\rho}^{2}  \, ( 1 + \rho ) \, e^{ - \beta \rho } \, d \rho 
		 + 
		\| a \|_{\infty}^{2} \int_{ D \setminus D( 4T )} \int_{ B( z, 2T ) } F_{t,t', d( z,z' )} \sum_{\gamma \in \Gamma}^{} F_{t,t', d( \gamma z, z' )} dz' dz,
	\end{align}
	where $ \beta = \beta( \lambda_{1}( X ) ) $, $ \lambda_{1}( X ) $ is the spectral gap of the Laplace--Beltrami operator $ -\Delta_{X} $, and
\begin{align*}
	\beta( x ) = 
	\begin{cases}	
		1 - \sqrt{ 1 - 4x }, & \text{ if } x \leq \frac{1}{4},\\
		1, & \text{ if } x > \frac{1}{4}.
	\end{cases}
\end{align*}
\end{lemma}

\begin{lemma}[]
	\label{Lmm: Small injectivity radius}
	Let $ X = \Gamma \setminus \mathbb{H} $ be a compact connected hyperbolic surface with fundamental domain $ D $, and let $ D( 4T ) $ denote the points with injectivity radius larger than $ 4T $. Then
	\begin{align*}
		&
		\int_{D \setminus D( 4T )} \int_{ B( z, 2T ) }   F_{t,t', d( z,z' )}  \sum_{\gamma \in \Gamma}^{} F_{t,t', d( \gamma z, z' )} dz' dz 
		\leq 4 \pi^{2} \, | D \setminus D( 4T ) | \, \frac{e^{4T}}{\operatorname{InjRad}_{X}} \int_{0}^{2T} \sinh( \rho ) F_{t,t',\rho}^{2} d\rho.
	\end{align*}
\end{lemma}

\begin{lemma}[]
	\label{Lmm: Weight integral}
	Let $ \beta $ be a positive constant smaller than $ 1 $ and let $ 0 < t' < t $. Then
	\begin{align*}
		&
		\int_{0}^{2T} \sinh( \rho ) \, ( 1 + \rho ) \, e^{ - \beta  \rho } F_{t,t',\rho}^{2} d \rho \leq \frac{C_{6}}{ \beta^{2} }  \, e^{ - \frac{\beta}{4}( t - t' ) },
	\end{align*}
	when
	\begin{align*}
		F_{t,t',\rho} \leq 50 \frac{1}{\sqrt{ \sinh( \max( t-t', \rho ) ) }} \boldsymbol 1_{ \rho \in ( 0, t + t' ) }, 
	\end{align*}
	where $C_{6} = 10000$.
\end{lemma}

\begin{lemma}[]
	\label{Lmm: Points do not matter}
	Let $ z,z', w, w' \in \mathbb{H} $ such that $ d( z,z' ) $ equals $ d( w,w' ) $. Then
	\begin{align}
		&
	     \nonumber
		\int_{ B( z,t ) \cap B( z', t' ) } \frac{dx}{ \sqrt{ \cosh( t ) - \cosh( d( x,z ) )  } \, \sqrt{ \cosh( t' ) - \cosh( d( x,z' ) ) } }	
	     \\ & = \ 
		\label{Expr: z z' integral}
		\int_{ B( w,t ) \cap B( w', t' ) } \frac{dx}{ \sqrt{ \cosh( t ) - \cosh( d( x,w ) )  } \, \sqrt{ \cosh( t' ) - \cosh( d( x,w' ) ) } }.
	\end{align}
		
\end{lemma}

\begin{proposition}[Geometric data]
\label{Prp: Geometric data}

Let $P_t$ and $P_{t,\tau}$ be integral operators on a compact connected hyperbolic surface $X = \Gamma \setminus \mathbb{H}$ with Laplace--Beltrami operator $ -\Delta_{X} $ such that $ P_{t} $ has integral kernel $ K_{t}^{\Gamma} $, defined as
\[
K_t^\Gamma(x,y) = \frac{1}{2 \sqrt{2} \pi} \sum_{\gamma \in \Gamma} K_t(x, \gamma y), \quad
K_t(x,y) = \frac{\mathbf{1}_{t > d(x,y)}}{\sqrt{\cosh(t) - \cosh(d(x,y))}},
\]
and $ P_{t,\tau} $ has integral kernel $K_{t,\tau}^{\Gamma} = \cos(t \tau) \, K_t^\Gamma$. Let $ a $ be a function on $ X $ with mean zero. Then
\[
	\left\| \frac{1}{T} \int_0^T P_{t,\tau}^{*} a P_t \, dt \right\|_{HS}^{2} \leq \frac{C_{7}}{T \, \beta^{3}} \, \left( \| a \|_{2}^{2} \, + \,  \| a \|_{\infty}^{2} \, | X \setminus X( 4T ) | \frac{e^{6 T}}{ \operatorname{InjRad}_{X} } \right)
\]
where $ \beta = \beta( \lambda_{1}( X ) ) $,  $|X \setminus X( 4T )|$  is the volume of the points of $ X $ having injectivity radius less than $ 4T $, $ \lambda_{1}( X ) $ is the spectral gap of $ -\Delta_{X} $, and where
\begin{align*}
	\beta( x ) = 
	\begin{cases}	
		1 - \sqrt{ 1 - 4x }, \qquad & \text{ if } x \leq \frac{1}{4},
		\\
		1, \qquad & \text{ if } x > \frac{1}{4}.
	\end{cases}
\end{align*}
The constant appearing is $C_{7} = 80000$.

\end{proposition}

\begin{proof}[Proof of Proposition \ref{Prp: Geometric data}]
	
Writing the Hilbert-Schmidt norm as a double integral over the fundamental domain $ D $ of $ X $ yields the following expression:
\begin{align*}
	&
	\int_{D} \int_{D} \left|  \frac{1}{T} \int_{0}^{T} \int_{D} K_{t,\tau}^{\Gamma} a( z ) K_{t}^{\Gamma}( z,y ) dt \right|^{2} dy dx.
\end{align*}
We also used that $ P_{t,\tau} $ is self-adjoint. Using that $ K^{\Gamma}_{t,\tau} = \cos(t \tau) K_t^\Gamma $ and by expanding the squares, we can write the previous expression as follows:
\begin{align*}
	&
	\int_{D \times D} \frac{1}{T^{2}} \int_{0}^{T} \int_{0}^{T} \cos( t \tau ) \cos( t' \tau ) \int_{D \times D} K_{t}^{\Gamma}( x,z ) a( z ) K_{t}^{\Gamma}( z,y ) \, K_{t'}^{\Gamma}( x,z' ) a( z' )^{*} K_{t'}^{\Gamma}( z', y )  dz' dz dt' dt dy dx.
\end{align*}
Using Fubini's theorem to reorder the integrals yields the following equal expression:
\begin{align*}
	&
	\frac{1}{T^{2}} \int_{0}^{T} \int_{0}^{T} \cos( t \tau ) \cos( t' \tau ) \int_{D \times D \times D \times D } K_{t}^{\Gamma}( x,z ) a( z ) K_{t}^{\Gamma}( z,y ) K_{t'}^{\Gamma}( x,z' ) a( z' )^{*} K_{t'}^{\Gamma}( z', y ) dx dy dz' dz dt' dt. 	
\end{align*}
This expression is bounded by adding the absolute values around it. Then we pull the absolute values inside to get the following upper bound:
\begin{align*}
	&
	\frac{1}{T^{2}} \int_{0}^{T} \int_{0}^{T} \left| \int_{D} \int_{D} \int_{D}^{} \int_{D} K_{t}^{\Gamma}( x,z ) a( z ) K^{\Gamma}_{t}( z,y ) \, K_{t'}^{\Gamma}( x,z' ) a( z' )^{*} K_{t'}^{\Gamma}( z', y )  dy dx dz' dz \right| dt' dt.
\end{align*}
The cosines were bounded from above by $ 1 $ in the previous step. The previous expression can be written in a more compact form when gathering the factors containing $ x $ into one integral and the factors containing $ y $ into another integral and noticing that these two integrals are the same, thus resulting in an integral squared:
\begin{align*}
	&
	\frac{1}{T^{2}} \int_{0}^{T} \int_{0}^{T} \left| \int_{D} \int_{D} a( z ) a( z' )^{*} \left( \int_{D} K_{t}^{\Gamma}( x,z ) K_{t'}^{\Gamma}( x,z' ) dx \right)^{2} dz' dz \right| dt' dt.
\end{align*}

Lemma \ref{Lmm: Eliminate sum 1} allows us to write the previous expression using the hyperbolic plane kernels $ K_{t} $ and $ K_{t'} $ instead of the surface kernels $ K_{t}^{\Gamma} $ and $ K_{t'}^{\Gamma} $: 
\begin{align*}
	&
	\frac{1}{8 \pi^{2} T^{2}} \int_{0}^{T} \int_{0}^{T} \left| \int_{D} \int_{D} a( z ) a( z' )^{*} \left( \int_{\mathbb{H}} \sum_{\gamma \in \Gamma}^{} K_{t}( x,z ) K_{t'}( x, \gamma z' ) dx \right)^{2} dz' dz \right| dt' dt,
\end{align*}
where the coefficient $ ( 8 \, \pi^{2} )^{-1} $ comes from the definition of $ K_{t}^{\Gamma} $:
\begin{align*}
	&
	K_{t}^{\Gamma}( x,y ) = \frac{1}{2 \sqrt{2 }\pi} \sum_{\gamma \in \Gamma}^{} K_{t}( x, \gamma y ).
\end{align*}


Lemma \ref{Lmm: Eliminate sum 2} allows us to rewrite the integral inside the absolute value in a convenient form, eliminating all but one sum. Thus the previous expression equals:
\begin{align}
	&
	\label{Expr: Introducing F}
	\frac{1}{ 8 \pi^{2} T^{2}} \int_{0}^{T} \int_{0}^{T} \left| \int_{D} a( z ) \int_{\mathbb{H}} a( z' )^{*}   F_{t,t', d( z,z' )}  \sum_{\gamma \in \Gamma}^{} F_{t,t', d( \gamma z, z' )} dz' dz \right| dt' dt,
\end{align}
where $ F_{t,t', \rho} $ is defined as
\begin{align*}
\int_{\mathbb{H}} K_{t}( x,w ) K_{t'}( x,w' ) dx,
\end{align*}
with $ w, w' \in \mathbb{H} $ separated by distance $ \rho $; Lemma \ref{Lmm: Points do not matter} ensures that $ F_{t,t',\rho} $ is well-defined. 

The definition of $ F_{t,t', d( z,z' )} $ shows that $ F_{t,t',d( z,z' )} = 0 $ if $ d( z,z' ) > 2T $, which follows from the definitions of $ K_{t} $ and $ K_{t'} $. Using this, Expression \eqref{Expr: Introducing F} can be rewritten by restricting the integration domain of \(z'\) as follows:
\begin{align}
	&
	\label{Expr: Before large small InjRad split}
	\frac{1}{8 \pi^{2} T^{2}} \int_{0}^{T} \int_{0}^{T} \left| \int_{D} a( z ) \int_{ B( z, 2T ) } a( z' )^{*}   F_{t,t', d( z,z' )}  \sum_{\gamma \in \Gamma}^{} F_{t,t', d( \gamma z, z' )} dz' dz \right| dt' dt.
\end{align}
We split the integral over $ D $ into points with injectivity radius greater than $ 4T $, $ D( 4T ) $, and points with injectivity radius less than $ 4T $, $ D \setminus D( 4T ) $:
\begin{align*}
	&
	\left| \int_{D} \cdots \right| \leq
	\left| \int_{D( 4T )} \cdots \right| 
	+  \left| \int_{D \setminus D( 4T )} \cdots \right|.
\end{align*}
By moving the absolute value inside the latter term and bounding \( |a| \) by its supremum norm, the previous expression yields the following upper bound:
\begin{align}
	&
	\nonumber
	\left| \int_{D( 4T )} a( z ) \int_{ B( z, 2T ) } a( z' )^{*}   F_{t,t', d( z,z' )}  \sum_{\gamma \in \Gamma}^{} F_{t,t', d( \gamma z, z' )} dz' dz \right| 
	\\
	+ \ & 
	\| a \|_{\infty}^{2} \int_{D \setminus D( 4T )} \int_{ B( z, 2T ) }  F_{t,t', d( z,z' )}  \sum_{\gamma \in \Gamma}^{} F_{t,t', d( \gamma z, z' )} dz' dz,
\end{align}
where the large and small injectivity radius terms are bounded separately. 

The small injectivity radius term is bounded in Lemma \ref{Lmm: Large injectivity radius}. Using this estimate, we obtain the following upper bound:
\begin{align*}
	&
	30 \| a \|_{2}^{2} \, \int_{0}^{2T}  \sinh( \rho ) \, F_{t,t',\rho}^{2} \, ( 1 + \rho ) \, e^{ -\beta \rho } d \rho 
	+ 2 \, \| a \|_{\infty}^{2} \int_{D \setminus D( 4T )} \int_{ B( z, 2T ) }   F_{t,t', d( z,z' )}  \sum_{\gamma \in \Gamma}^{} F_{t,t', d( \gamma z, z' )} dz' dz.
\end{align*}
The second term of the previous expression is bounded in Lemma \ref{Lmm: Small injectivity radius}, yielding the following upper bound:
\begin{align*}
	&
	30 \| a \|_{2}^{2} \, \int_{0}^{2T}  \sinh( \rho ) \, F_{t,t',\rho}^{2} \, ( 1 + \rho ) \, e^{ -\beta  \rho } d \rho 
	\ + \ 8\pi^{2}  \, \| a \|_{\infty}^{2} | D \setminus D( 4T ) | \frac{e^{4T}}{\operatorname{InjRad}_{X}} \int_{0}^{2T} \sinh( \rho ) F_{t,t',\rho}^{2} d\rho.
\end{align*}
This is further bounded by
\begin{align*}
	&
	8 \pi^{2}\left( \| a \|_{2}^{2} \, + \, \| a \|_{\infty}^{2} \, | D \setminus D( 4T ) | \, \frac{ e^{ 6T }}{ \operatorname{InjRad}_{X} } \right) \int_{0}^{2T} \sinh( \rho ) e^{ -\beta \rho} ( 1 + \rho ) F_{t,t',\rho}^{2} d \rho,
\end{align*}
using that $ 8 \pi^{2} > 30 $ and $ e^{ 2T } \, e^{ -\beta \rho } \, ( 1 + \rho ) \geq 1 $.

Hence, Expression \eqref{Expr: Introducing F} is bounded by
\begin{align*}
	&
	\frac{1}{ T^{2}}  \,  \left( \| a \|_{2}^{2} \, + \,  \| a \|_{\infty}^{2} \, | D \setminus D( 4T ) | \frac{e^{6T}}{ \operatorname{InjRad}_{X} } \right)\int_{0}^{T} \int_{0}^{T} \int_{0}^{2T} \sinh( \rho ) e^{-\beta \rho} ( 1 + \rho ) F_{t,t', \rho}^{2} d\rho dt' dt.
\end{align*} 
By symmetry of $ F_{t,t',\rho} $ in $ t $ and $ t' $, we can restrict the inner integral to $ 0 \le t' \le t $ and multiply by $ 2 $:
\begin{align*}
	&
	\frac{2}{ T^{2}} \, \left( \| a \|_{2}^{2} \, + \, \| a \|_{\infty}^{2} \, | D \setminus D( 4T ) | \frac{e^{6T }}{ \operatorname{InjRad}_{X} } \right)\int_{0}^{T} \int_{0}^{t} \int_{0}^{2T} \sinh( \rho ) e^{ -\beta \rho } ( 1 + \rho ) F_{t,t', \rho}^{2} d\rho dt' dt.
\end{align*}
Applying Proposition \ref{Lmm: Weight integral} to the innermost integral, we obtain the upper bound
\begin{align*}
\frac{2}{T^{2}} \left( \| a \|_{2}^{2} + \| a \|_{\infty}^{2} \, | D \setminus D( 4T ) | \frac{e^{6T}}{\operatorname{InjRad}_{X}} \right)
\int_{0}^{T} \int_{0}^{t} \frac{C_{6}}{\beta^{2}} \, e^{- \frac{\beta}{4}(t-t')} \, dt' \, dt,
\end{align*}
where \(C_{6} = 10000\).
Integrating in $ t' $ and $ t $ gives an upper bound of $ 4T C_{6} / \beta^{3} $, so the entire expression is bounded by
\begin{align*}
	&
	\frac{C_{7}}{T \, \beta^{3}} \, \left( \| a \|_{2}^{2} \, + \, \| a \|_{\infty}^{2} \, | D \setminus D( 4T ) | \frac{e^{6 T}}{ \operatorname{InjRad}_{X} } \right),
\end{align*}
with $2 \cdot 4 \cdot C_{6} = 80000 = C_{7}$.

\end{proof}

\end{document}
